\documentclass[conference]{IEEEtran}
\IEEEoverridecommandlockouts
% The preceding line is only needed to identify funding in the first footnote. If that is unneeded, please comment it out.
\usepackage{cite}
\usepackage{amsmath,amssymb,amsfonts}
\usepackage{algorithmic}
\usepackage{graphicx}
\usepackage{textcomp}
\usepackage{xcolor}
\usepackage[colorlinks=true, linkcolor=blue, citecolor=blue, urlcolor=blue]{hyperref}
\def\BibTeX{{\rm B\kern-.05em{\sc i\kern-.025em b}\kern-.08em
    T\kern-.1667em\lower.7ex\hbox{E}\kern-.125emX}}
\begin{document}

\title{CodeSentinel: A Privacy-Preserving Deep Intelligence Platform for Semantic Security Auditing and Dynamic Sandboxing\\
{\footnotesize \textsuperscript{*}Note: Sub-titles are not captured in Overleaf by default if specific templates are used.}
}

\author{\IEEEauthorblockN{1\textsuperscript{st} Ali Hamza}
\IEEEauthorblockA{\textit{Research Fellow} \\
\textit{CodeSentinel AI}\\
City, Country \\
email@address.com}
}

\maketitle

\begin{abstract}
CodeSentinel represents a paradigm shift in automated security auditing by integrating local Large Language Models (LLMs) with containerized dynamic execution environments. This paper details the architectural design and performance metrics of CodeSentinel...
\end{abstract}

\begin{IEEEkeywords}
Cybersecurity, Artificial Intelligence, LLM, Static Analysis, Dynamic Analysis, Docker, Privacy-Preserving AI.
\end{IEEEkeywords}

\section{Introduction}
% The introduction will be synthesized based on your input.

\section{Architecture Overview}
% The architecture details from the guide will be placed here.

\section{Implementation}
% Technical implementation details.

\section{Analysis and Results}
% Performance HUD metrics and findings.

\section{Conclusion}
% Summary of contributions.

\bibliographystyle{IEEEtran}
\bibliography{references}

\end{document}
